\section{Open Weights and Audio}
\sectionkicker{OPEN WEIGHTS}
\begin{wideflow}
{\Large\bfseries 公開の言明と取得の実態で仕分ける――重み・音声7件\par}
\smallskip
{\color{SurveyMuted}ダウンロードの可否とライセンス確認の有無を軸に、7件の重み・音声提供を一覧化する。\par}
\end{wideflow}

\subsection*{公開の言明と取得の実態で仕分ける――重み・音声7件}

今週のオープンウェイト周辺は、規模の競争ではなく公開状態の差で読むべき週だ。Aleph AlphaのKolibriは10月3日に78B／3BのMoE構成で公開され、コンテキスト100万token、ライセンスApache-2.0、HFのAleph-Alpha／Kolibri-1で取得可能とされる。学習20T token、計算資源768基のB200という条件も添えられている。いずれも発行元の公開ページの記述であり、配布場所まで含めた具体性がある。\cite{kolibri}

性能面ではAIME 2026で96.0、GPQAで84.3といったベンダー数値が示されている。これらは発行元が自社条件で測った申告値で、独立再現はない。ただし公開ページの表は途中で取得が途切れており、表の残部は未確認である。規模と条件の層では語れるが、評価の全容は留保付きになる。\cite{kolibri}

小型の決定モデルでは、Liquid AIのd1-3Bとd1-omni-600Mが、2026年10月7日の発表で取得可能とされた。Decision Indexは48.57と15.95、RTX4090上で8ミリ秒というベンダー数値が添えられている。600M級のオムニモデルと8ミリ秒という遅延の組み合わせは、端末側での実行を考える開発者には具体的な検討材料だ。ただしIndexの運営方法や測定条件は不明で、価格とライセンスのSPDXも取得本文にない。\cite{liquidd1}

Kolibriの大型MoEとd1の小型決定モデルは、対照的な二極として読める。片方は100万token級の文脈を抱える巨体、もう片方はミリ秒応答の小体だ。Decision Indexのような独自指標は、異なる規模のモデルを決定タスクの軸で比べる試みだが、指標自体の方法が不明なため、順位の確定には使えない。規模の提示としては有益、評価の確定としては未熟という二層で受け止めたい。\cite{liquidd1}

\subsection*{埋め込みの更新――EmbeddingGemma 2の仕様箱}

Googleは10月6日のブログでEmbeddingGemma 2を公開した。全体740M parametersの構成で、内訳はテキスト専用270M、画像用170M、音声用300Mの3エンコーダからなる。テキスト、コード、画像、映像、音声を共通の埋め込み空間で扱う設計で、768次元のMRL（Matryoshka）表現を持ち、128・256・512次元への短縮が可能とされる。コンテキストは8000 tokenで、音声は最大5.5分、画像は最大29枚、映像は最大58フレームが目安として示されている。ライセンスはApache-2.0、取得場所はHFのgoogle／embeddinggemma-2とKaggleとされる。MRL短縮の実用的な意味は、検索インデックスの容量や計算負荷の調整にある。短い次元で候補を絞り、長い次元で再評価するといった使い分けが考えられるが、これは実務上の応用例であり、開発元がその構成の優位性を実証したという意味ではない。次元数とMRL対応、コンテキスト長の三点はそのまま設計条件になる。\cite{embedgemma2}

検索・コード検索の指標としてMTEB-Codeで＋9.92というベンダー数値が出ているが、全指標はモデルカード側に委ねられ、そのカード本体は取得タイムアウトで未取得だ。アーキテクチャ表も断片しかなく、学習や評価の方法は分からない。仕様箱としての記事は成立するが、検索リーダーとしての優劣判断はカード取得待ちになる。\cite{embedgemma2}

\subsection*{比較対象――まだ取得が確定しない二件}

ReflectionのBeamは10月5日発表（10月8日更新）の501B／23B MoEで、重み・報告・カードは発表時点では2026年10月後半の公開予定とされ、ライセンスはApache-2.0の意向、現状は早期アクセス登録のみだ。ベンダー数値としてDeepSWE 44.4、TB2.1で80.1、AIMEで97.8、学習23.8T token、GB300級でのRL規模が添えられている。規模の提示は大きいが、取得が確定したものは何もない段階であり、Kolibriとの対比で「公開済み対予告」の構図を作る比較対象にとどめる。\cite{beam}

Mistral Large 4は10月6日のブログとdocsで、有効520億・総計1～1.05兆、コンテキスト100万、プレビューAPIは発表日から提供、重みは2026年10月末、セール価格は入力・出力100万token当たり\$0.68／\$2.09とされる。APIとしては発表日から試せる段階だが、重みとしては2026年10月末待ちだ。ここで二つの区別を守りたい。第一に、有効520億と総計1.05兆は別の量であり、前者が1 tokenごとの演算に寄与する規模、後者が保持メモリの目安になる。第二に、一部の二次記録には49B activeという数値があるが、対応する原X投稿は直接確認できていない。Mistral公式docsでは52B activeと記載されているため、本文では52Bを一次資料に基づく値として扱う。二つの数値を同等の確実性で並記はしない。公式ピンは52Bという区別を崩さず、裁定はしない。\cite{mistral4}

\subsection*{音声の二件――Clef-omniとNova Sonic}

CloudflareのClef-omniは、公開日時の確認が明確な好例でもある。Clef-omniの公開時刻は、発行元ブログに記録された公開日時のメタデータに基づく。2026年10月9日18時27分34秒（UTC）の記録があり、SNS上の後続投稿日時とは別の扱いになる。SNS上の後続投稿日時は、モデルの公開日時を直接示すものではない。\cite{clefomni}

仕様としてはQwen3-Omni-30B-A3BのMoEで、omni料金\$0.15／M、flashの値下げ\$0.09→\$0.038、HFでの重み公開、ベンダーベンチが添えられている。ただしHFのダウンロード実績とライセンスSPDXは未検証で、発行者の公開言明と実際に取得されたファイルは別物だ。料金と仕様の層では語れるが、取得の層では未確認という留保が付く。\cite{clefomni}

AWSのNova 2.5 Sonicは10月5日に音声対音声で一般提供（GA）となった。Bedrockの4リージョン（米東・米西・北欧・東京含むアジア太平洋北東）で利用でき、コンテキスト25.6万、前世代と同等料金とされる。GAはプレビューではなく本提供の段階を指し、対象地域の確定が今回の到達内容だ。ただし料金の数値は本文になく、ベンチマークやライセンス条項もない。価格と品質の判断は材料待ちであり、GAの確定と支払い条件の確定は別の事柄として扱う。\cite{novasonic}

\begin{table}[htbp]
\centering\footnotesize
\begin{tabularx}{0.98\linewidth}{@{}lX@{}}
\toprule
項目 & 内容 \\
\midrule
Kolibri & ライセンス表記：Apache-2.0（言明）；配布場所：HF Aleph-Alpha／Kolibri-1（言明）；実ファイルの取得確認：言明場所の提示のみ、SPDX・実ダウンロード未検証；状態：公開言明あり（表残部未取得） \\
Liquid d1 & ライセンス表記：SPDXなし；配布場所：HF LiquidAI/d1-3b等（言明）；実ファイルの取得確認：未検証；状態：取得可能との言明あり \\
EmbeddingGemma 2 & ライセンス表記：Apache-2.0（言明）；配布場所：HF＋Kaggle（言明）；実ファイルの取得確認：カード未取得；状態：公開言明あり、詳細はカード待ち \\
Beam & ライセンス表記：Apache-2.0意向；配布場所：なし；実ファイルの取得確認：不可；状態：発表時点で2026年10月後半の公開予定、早期登録のみ \\
Mistral Large 4 & ライセンス表記：WEIGHTSタブ未取得；配布場所：プレビューAPIあり；実ファイルの取得確認：不可；状態：APIは発表日から提供、重みは2026年10月末 \\
Clef-omni & ライセンス表記：SPDXなし；配布場所：HF（言明・リンク）；実ファイルの取得確認：未検証；状態：公開言明あり、時刻確定済み \\
Nova Sonic & ライセンス表記：本文なし；配布場所：Bedrock 4リージョン；実ファイルの取得確認：対象外（API提供）；状態：GA、料金数値なし \\
\bottomrule
\end{tabularx}
\caption{比較表：重み・音声7件の公開状態}
\end{table}
\noindent \cite{kolibri}\cite{liquidd1}\cite{embedgemma2}\cite{beam}\cite{mistral4}\cite{clefomni}\cite{novasonic}

この表の読み方が本記事の核心だ。発行者の公開表明、配布先の発行元による提示、その提示先の実在確認、重みファイルの取得成功は四つの別の段階であり、今回確認した資料から裏付けられるのは、発行元による公開表明と配布先の提示までである。たとえばHFのリポジトリ名が示されていても、その提示先の実在確認と重みファイルの取得がなければ、利用許諾の判断は下せない。試用はできても商用組み込みの判断は時期尚早ということになる。\cite{kolibri}\cite{beam}\cite{mistral4}\cite{clefomni}

規模の読み方も整理する。MoEのtotalとactiveは別の量だ。各モデルのtotal／activeパラメータ数は、Kolibriが約78.1B／約3.46B、Beamが501B／23B、Mistral Large 4が1.05T／52Bである。totalはモデル全体の重み規模で、保持に必要なメモリの目安になる。activeはtoken処理時に主として利用されるパラメータ規模で、推論時の演算コストの目安になる。ただし、active数だけで実際のFLOPsや処理速度が確定するわけではない。量子化の有無、実装、ルーティング方式、ハードウェアの違いで変わるため、規模表とベンチ表を別々の表として扱い、混ぜて総合ランキングを作らないことが正確さの条件になる。\cite{kolibri}\cite{beam}\cite{mistral4}

配布場所の確認手順も示しておく。HFのリポジトリ名が示されていても、ライセンス条項ファイルの有無と重みファイルの一覧を確認するまでは、利用許諾の判断は下せない。Clef-omniのようにリンクが示されていても、実ファイルの取得確認がなければ言明の段階にとどまる。確認の四段階――発行元による公開表明、発行元による配布先の提示、その提示先の実在確認、重みファイルの取得成功――を踏むことが、オープンウェイト利用の基本動作になる。ライセンス条項の確認は、単なる配布先提示やダウンロード成功とは別に必要だ。\cite{kolibri}\cite{clefomni}

音声二件の対比も述べておく。NovaはサービスとしてのGA、Clefは重み公開言明付きのオムニ発表で、提供形態が異なる。料金はNovaが数値なし、Clefが言明値ありと逆転しているが、いずれも確定した支払い条件ではない。対象地域の確定（Novaの4リージョン）とプレビュー段階の区別（MistralのAPI先行）も含め、提供形態の確認から始めるのが順序だ。\cite{novasonic}\cite{clefomni}\cite{mistral4}\cite{beam}

日本の開発者への要点は、取得可能とされるのはKolibriとd1、仕様確認の段階なのはEmbeddingGemma、音声はNovaのGAとClefの言明という仕分けだ。重みの公開予定は分かれており、Beamは発表時点で2026年10月後半の公開予定、Mistral Large 4は発表時点で2026年10月末の公開予定とされた。発表予定を実際の公開完了と同一視してはならない。いずれも発行元の表明に基づく仕分けであり、実ファイルの取得確認はこれからだ。表明の段階で検証を進め、取得の段階で利用判断に移る。この順序づけが今週の実利になる。\cite{kolibri}\cite{liquidd1}\cite{embedgemma2}\cite{novasonic}\cite{clefomni}\cite{beam}\cite{mistral4}

\begin{technicalnote}[Note]
公開表明・配布場所・実ファイル取得は別の段階として扱った。価格・ベンチ数値は発行元申告であり、独立再現なし。
\end{technicalnote}
